\section*{Appendix S1. Gauge-Orbit Diagnostic Decomposition}
\label{supp:orbit-diagnostic}
\renewcommand{\theequation}{S1.\arabic{equation}}
\setcounter{equation}{0}

\subsection*{S1.1 Purpose and scope}

This appendix defines the invariant and frame-dependent parts of the operator
comparison used in the main paper. The decomposition applies to the diagnostic,
not to the Hamiltonian as a physical operator. It is defined relative to a
frozen comparison contract; changing the retained space, admissible alignment
family, weights, normalization, or finite translation domain changes the
diagnostic.

The construction is analogous in purpose to the separation of the Wannier
spread into gauge-invariant and gauge-dependent contributions
\cite{marzari1997,marzari2012,panati2013}. It is not the Wannier spread
functional and does not inherit that functional's physical or variational
interpretation.

\subsection*{S1.2 Comparison contract}

Let $V_{\mathrm{ref}}$ and $V_{\mathrm{red}}$ be retained finite-dimensional
state spaces of equal rank $r$. The operator channel requires a nonempty,
declared family
\[
 \mathcal U\subseteq
 \left\{\bm C:V_{\mathrm{red}}\rightarrow V_{\mathrm{ref}}
 \mid \bm C^\dagger\bm C=\bm C\bm C^\dagger=\bm I_r\right\}
\]
of admissible unitary identification maps. Restrictions imposed by continuity,
reciprocal-boundary sewing, symmetry, orbital semantics, or locality belong to
the definition of $\mathcal U$ rather than to a later interpretation of the
residual.

For a finite translation domain $\mathcal S_H$, nonnegative weights
$\omega_{\bm R}$, and a positive frozen normalization
\[
 Z_{\mathrm{ref}}=
 \sum_{\bm R\in\mathcal S_H}\omega_{\bm R}
 \left\|\bm H_{\mathrm{ref}}(\bm R)\right\|_{\mathrm F}^2,
\]
the aligned operator loss is
\begin{equation}
 \mathcal L_H(\bm C,\bm\theta)=
 \frac{1}{Z_{\mathrm{ref}}}
 \sum_{\bm R\in\mathcal S_H}\omega_{\bm R}
 \left\|
 \bm H_{\mathrm{ref}}(\bm R)
 -\bm C\bm H_{\mathrm{red}}(\bm R;\bm\theta)\bm C^\dagger
 \right\|_{\mathrm F}^2.
 \label{eq:supp-aligned-loss}
\end{equation}
A zero or unspecified normalization does not define this diagnostic. A
different-rank comparison likewise requires a separately declared projection
or embedding; its information loss cannot be absorbed into
Eq.~\eqref{eq:supp-aligned-loss}.

\subsection*{S1.3 Orbit loss and frame excess}

\paragraph{Attainment under a compact contract.}
For fixed $\bm\theta$, the finite-sum loss in
Eq.~\eqref{eq:supp-aligned-loss} is a continuous real-valued function of
$\bm C$. If $\mathcal U$ is nonempty and compact, the extreme-value theorem
therefore supplies $\bm C_\star\in\mathcal U$ such that
\begin{equation*}
 \inf_{\bm C\in\mathcal U}\mathcal L_H(\bm C,\bm\theta)
 =\min_{\bm C\in\mathcal U}\mathcal L_H(\bm C,\bm\theta)
 =\mathcal L_H(\bm C_\star,\bm\theta).
\end{equation*}
In the finite-rank contract used here, $U(r)$ is compact, so it is sufficient
for the admissible family to be a nonempty closed subset of $U(r)$. Sewing,
symmetry, site, and orbital restrictions imply this conclusion only when their
frozen constraint maps define a closed set; the names of those restrictions do
not by themselves prove compactness. For a merely nonempty noncompact family,
the nonnegative loss still has a well-defined infimum, but a minimizing
alignment need not exist.

For nonempty $\mathcal U$, define the contract-invariant orbit loss
\begin{equation}
 \mathcal L_{H,\mathrm{orb}}(\bm\theta)
 =\inf_{\bm C\in\mathcal U}\mathcal L_H(\bm C,\bm\theta).
 \label{eq:supp-orbit-loss}
\end{equation}
For a particular declared alignment $\bm C_0\in\mathcal U$, define the
frame-dependent excess
\begin{equation}
 \mathcal L_{H,\mathrm{frame}}(\bm C_0,\bm\theta)
 =\mathcal L_H(\bm C_0,\bm\theta)
 -\mathcal L_{H,\mathrm{orb}}(\bm\theta).
 \label{eq:supp-frame-excess}
\end{equation}
For every $\bm C_0\in\mathcal U$, the defining greatest-lower-bound
property gives
$\mathcal L_{H,\mathrm{orb}}(\bm\theta)\leq
\mathcal L_H(\bm C_0,\bm\theta)$. Subtraction therefore proves
\begin{equation}
 \mathcal L_H(\bm C_0,\bm\theta)
 =\mathcal L_{H,\mathrm{orb}}(\bm\theta)
 +\mathcal L_{H,\mathrm{frame}}(\bm C_0,\bm\theta),
 \qquad
 \mathcal L_{H,\mathrm{frame}}\geq 0.
 \label{eq:supp-loss-decomposition}
\end{equation}
This is an exact scalar decomposition of the diagnostic. It is not an
orthogonal decomposition of the residual matrix.

The orbit term is the greatest lower bound of the operator discrepancy over the
frozen admissible family. The frame term measures the excess associated with
the selected admissible representative. Both statements are relative to
$\mathcal U$: enlarging the family can decrease the orbit loss, while adding
physical restrictions can increase it.

\subsection*{S1.4 Frame invariance}

Consider independent frame changes $\bm A$ on $V_{\mathrm{ref}}$ and $\bm B$
on $V_{\mathrm{red}}$. They transform the represented operators and alignment
map as
\[
 \bm H_{\mathrm{ref}}' =\bm A\bm H_{\mathrm{ref}}\bm A^\dagger,
 \qquad
 \bm H_{\mathrm{red}}' =\bm B\bm H_{\mathrm{red}}\bm B^\dagger,
 \qquad
 \bm C'=\bm A\bm C\bm B^\dagger.
\]
If the admissible family is transported bijectively to
$\mathcal U'=\{\bm A\bm C\bm B^\dagger:\bm C\in\mathcal U\}$, then the
residual corresponding to $\bm C'=\bm A\bm C\bm B^\dagger$ is $\bm A$ times
the original residual times $\bm A^\dagger$. Unitary invariance of the
Frobenius norm gives
\[
 \mathcal L_H'(\bm C',\bm\theta)=\mathcal L_H(\bm C,\bm\theta).
\]
Because $\bm C\mapsto\bm A\bm C\bm B^\dagger$ is a bijection,
\begin{align*}
 \inf_{\bm C'\in\mathcal U'}\mathcal L_H'(\bm C',\bm\theta)
 &=\inf_{\bm C\in\mathcal U}
   \mathcal L_H'(\bm A\bm C\bm B^\dagger,\bm\theta)\\
 &=\inf_{\bm C\in\mathcal U}\mathcal L_H(\bm C,\bm\theta).
\end{align*}
Thus $\mathcal L_{H,\mathrm{orb}}$ is unchanged. Applying the same identity to
the corresponding representative
$\bm C_0'=\bm A\bm C_0\bm B^\dagger$ and subtracting the equal orbit losses
also proves equality of the two frame excesses. The numerical matrices
representing $\bm C_0$ and the residual may nevertheless change.

This invariance is contractual rather than absolute. A transformation that
violates a frozen symmetry, site, orbital, or sewing condition does not belong
to the same admissible comparison problem.

\subsection*{S1.5 Outcomes when subtraction is unavailable}

The diagnostic has three distinct outcomes:
\begin{enumerate}
 \item If $\mathcal U$ is nonempty, the orbit loss is defined. A feasible
 alignment supplies an upper bound; a certified lower bound requires a separate
 global argument.
 \item If the retained ranks differ but a scientifically authorized projection
 or embedding is supplied, comparison may proceed in the declared target space.
 Projection or embedding error remains a separate diagnostic.
 \item If no admissible identification exists, the matrix-difference channel is
 undefined. The recorded outcome is a state-space mismatch, not a large
 operator residual.
\end{enumerate}

Gauge-invariant channels can remain available in the third case. Spectra and
trace invariants may be compared after their own indexing and normalization
contracts are fixed. Projector distances require both retained subspaces to be
embedded in the same parent Hilbert space. Wilson-loop eigenphase multisets can
be compared without choosing an internal frame, provided the loop, rank, gap,
and reciprocal sewing contracts agree. These invariant channels do not create
an operator identification map by themselves.

\subsection*{S1.6 Relation to the Wannier decomposition}

For a retained Bloch subspace, the Wannier spread admits the established
separation
\[
 \Omega=\Omega_{\mathrm I}+\widetilde\Omega,
\]
where $\Omega_{\mathrm I}$ depends only on the subspace projectors and
$\widetilde\Omega$ depends on the selected Bloch frame
\cite{marzari1997,marzari2012}. Equations~\eqref{eq:supp-orbit-loss}--
\eqref{eq:supp-loss-decomposition} use the same conceptual distinction between
an equivalence class and its representative, but they answer a different
question. They quantify cross-model operator discrepancy over a declared
admissible orbit; they do not quantify Wannier localization.

Accordingly, a small orbit loss establishes agreement only for the finite
operator diagnostic that was defined. A small frame excess places the achieved
loss near the orbit infimum; it does not imply that the alignment matrix is
close to a unique optimizer. Neither quantity, by itself, establishes material
validity, continuum convergence, or locality outside the frozen comparison
domain.
